\begin{statementbox}
\begin{center}\begin{tikzpicture}[font=\small,>=stealth,line join=round]
\draw[gray!40,rounded corners=3pt] (0,0) rectangle (14.2,8.5);\node at (7.1,7.9) {先看内层：合并同一个数，还是得到一对相反数？};
\node at (2.4,6.9) {$x^2+3$：内偶};\begin{scope}[shift={(2.4,3.8000000000000003)},x=1.1333333333333333cm,y=0.4cm]\draw[->,gray] (-1.5,0)--(1.5,0) node[right]{$x$};\draw[->,gray] (0,-3)--(0,6) node[above]{$y$};\draw[blue,thick,solid,domain=-1.3:1.3,samples=160] plot(\x,{(\x)^2+3});\fill[orange!85!black] (-1,4) circle(2pt);\fill[orange!85!black] (1,4) circle(2pt);\end{scope}\node[align=center,text width=4cm] at (2.4,1.7) {$F(-1)=F(1)=4$};\node at (7.050000000000001,6.9) {$x^3$：内奇外奇};\begin{scope}[shift={(7.050000000000001,3.8000000000000003)},x=1.1333333333333333cm,y=0.4cm]\draw[->,gray] (-1.5,0)--(1.5,0) node[right]{$x$};\draw[->,gray] (0,-3)--(0,6) node[above]{$y$};\draw[violet,thick,solid,domain=-1.3:1.3,samples=160] plot(\x,{(\x)^3});\fill[orange!85!black] (-1,-1) circle(2pt);\fill[orange!85!black] (1,1) circle(2pt);\end{scope}\node[align=center,text width=4cm] at (7.050000000000001,1.7) {$F(-1)=-1,\ F(1)=1$};\node at (11.7,6.9) {$x^2$：内奇外偶};\begin{scope}[shift={(11.7,3.8000000000000003)},x=1.1333333333333333cm,y=0.4cm]\draw[->,gray] (-1.5,0)--(1.5,0) node[right]{$x$};\draw[->,gray] (0,-3)--(0,6) node[above]{$y$};\draw[orange!85!black,thick,solid,domain=-1.3:1.3,samples=160] plot(\x,{(\x)^2});\fill[orange!85!black] (-1,1) circle(2pt);\fill[orange!85!black] (1,1) circle(2pt);\end{scope}\node[align=center,text width=4cm] at (11.7,1.7) {$F(-1)=F(1)=1$};
\node[text width=12.6cm,align=center] at (7.1,.65) {三图同尺度；分别取 $g=x^2,f(u)=u+3$，及 $g=x,f(u)=u^3$ 或 $u^2$。};
\end{tikzpicture}\end{center}


\textbf{\textcolor{CStatement}{这条结论解决什么问题？}}
遇到 $F(x)=f(g(x))$ 这样的函数，如果整个式子很长，直接把 $-x$ 代入再化简可能很麻烦。我们想先看\textbf{里面那一步}，就能尽快判断整个图像是不是关于 $y$ 轴或原点对称。

\medskip
\textbf{\textcolor{CStatement}{先看懂三个字母}}
\[
\underbrace{x\ \xrightarrow{\ g\ }\ g(x)}_{\text{先算里面}}
\ \xrightarrow{\ f\ }\ \underbrace{f(g(x))}_{F(x)\text{，再算外面}}.
\]
例如 $F(x)=(x^2+1)^3$：先算 $g(x)=x^2+1$，再把算出的数立方，也就是 $f(u)=u^3$。这里的 $u$ 只是临时用来表示“里面算出来的数”。

\medskip
\textbf{\textcolor{CStatement}{使用前提}}
\begin{enumerate}[leftmargin=2.2em,itemsep=0.25em]
\item $x$ 可以取的范围必须左右对称：如果允许 $x=a$，也要允许 $x=-a$。
\item 对每个允许的 $x$，内层算出的 $g(x)$ 都必须能代入外层 $f$。\textbf{否则 $F(x)$ 根本没有定义。}
\end{enumerate}

\medskip
\end{statementbox}\clearpage\begin{statementbox}\textbf{\textcolor{CStatement}{核心结论：内偶则偶，内奇同外}}
\begin{center}\begin{tikzpicture}[font=\small,>=stealth,line join=round]
\draw[gray!40,rounded corners=3pt] (0,0) rectangle (14.2,8.3);\node at (7.1,7.700000000000001) {内偶的原因：第一步已经把正负输入合为同一个数};
\node at (7.1,6.9){取 $x=-2,2$，先平方，再加 $3$};\draw[->,gray] (1,5.7)--(12.8,5.7);\node[left] at (1,5.7){$x$};\foreach \X/\v in {2.4/-2,3.5/0,4.6/2}{\draw[gray] (\X,5.7)--(\X,5.6);\node[above=6pt] at (\X,5.7){$\v$};}\fill[blue] (2.4,5.7) circle(2.3pt);\fill[orange!85!black] (4.6,5.7) circle(2.3pt);\draw[->,gray] (1,3.6)--(12.8,3.6);\node[left] at (1,3.6){$u$};\draw[gray] (1.3,3.6)--(1.3,3.5);\node[above=6pt] at (1.3,3.6){$0$};\fill[violet] (3.5,3.6) circle(2.3pt);\node[below left=7pt] at (3.5,3.6){$4$（两路重合）};\draw[->,blue,dashed,shorten >=3pt] (2.4,5.7)--(3.5,3.6);\draw[->,orange!85!black,dashed,shorten >=3pt] (4.6,5.7)--(3.5,3.6);\draw[->,gray] (1,1.8)--(12.8,1.8);\node[left] at (1,1.8){$y$};\draw[gray] (1.3,1.8)--(1.3,1.7);\node[above=6pt] at (1.3,1.8){$0$};\fill[violet] (5.15,1.8) circle(2.3pt);\node[above right=7pt] at (5.15,1.8){$7$（两路仍重合）};\draw[->,violet,dashed,shorten <=3pt,shorten >=3pt] (3.5,3.6)--(5.15,1.8);\node[anchor=west] at (8,4.8){$g(x)=x^2$};\node[anchor=west] at (8,2.8){$f(u)=u+3$};
\node[text width=12.6cm,align=center] at (7.1,.65) {各行单位长度相同，零点位置可不同；所有标数均由实际运算得到。};
\end{tikzpicture}\end{center}


\begin{description}[style=nextline,leftmargin=2.5em,labelsep=0.4em,itemsep=0.6em]
\item[\textbf{第一种：内偶则偶。}]
如果 $g(-x)=g(x)$，那么不论外层 $f$ 是奇函数、偶函数，还是两者都不是，都有
\[
F(-x)=F(x),\qquad F\text{ 为偶函数}.
\]
\textbf{通俗解释：} 把 $x$ 换成 $-x$，里面已经算出了同一个数；外面再怎么算，结果都一样。
\item[\textbf{第二种：内奇、外奇，则整体奇。}]
如果 $g(-x)=-g(x)$，而且 $f(-u)=-f(u)$，那么
\[
F(-x)=-F(x),\qquad F\text{ 为奇函数}.
\]
\textbf{通俗解释：} 里面先把结果变成相反数，外面再把最后结果变成相反数。
\item[\textbf{第三种：内奇、外偶，则整体偶。}]
如果 $g(-x)=-g(x)$，而且 $f(-u)=f(u)$，那么
\[
F(-x)=F(x),\qquad F\text{ 为偶函数}.
\]
\textbf{通俗解释：} 虽然里面算出相反数，但外面会把两个相反数算成相同的结果。
\end{description}

\textbf{\textcolor{CStatement}{一句特别重要的限制}}
“内奇同外”只表示：\textbf{当外层明确是奇函数或偶函数时，可以把这个性质传给整体。} 如果外层既不是奇函数也不是偶函数，不能单凭口诀判断；也不能根据整体为偶函数，就反推里面一定是偶函数。

\medskip
\end{statementbox}\clearpage\begin{statementbox}\textbf{\textcolor{CStatement}{不能扩大使用条件：外偶未必整体偶}}
\begin{center}\begin{tikzpicture}[font=\small,>=stealth,line join=round]
\draw[gray!40,rounded corners=3pt] (0,0) rectangle (14.2,8.3);\node at (7.1,7.700000000000001) {外层平方为偶，整体却偏离了纵轴};
\begin{scope}[shift={(5.1,2.068181818181818)},x=1.3333333333333333cm,y=0.8363636363636363cm]\draw[->,gray] (-3,0)--(1.5,0) node[right]{$x$};\draw[->,gray] (0,-0.5)--(0,5) node[above]{$y$};\draw[blue,thick,solid,domain=-2.8:1.2,samples=180] plot(\x,{(\x+1)^2});\fill[orange!85!black] (-1,0) circle(2.2pt);\fill[orange!85!black] (1,4) circle(2.2pt);\end{scope}\node[anchor=west,align=left,text width=5.0cm] at (8.4,4.0) {$F=(x+1)^2$\\[10pt]内层 $g=x+1$ 非奇非偶\\[10pt]$F(1)=4,\ F(-1)=0$\\[10pt]对称轴 $x=-1$，不是 $y$ 轴};
\node[text width=12.6cm,align=center] at (7.1,.65) {外层只会把互为相反数的输入算成相同结果；这里的2和0并非相反数。};
\end{tikzpicture}\end{center}


\textbf{不能说：外层是偶函数，整体就一定是偶函数。} 比如 $f(u)=u^2$ 是偶函数，但取 $g(x)=x+1$，得到
\[
F(x)=(x+1)^2,\quad F(1)=4,\quad F(-1)=0.
\]
结果不同，所以 $F$ 不是偶函数；且 $F(-1)\neq -F(1)$，也不是奇函数。
\end{statementbox}
